\section{Methodology}
\label{sec:methodology}

\subsection{Primary Dataset and Split}

The primary experiments used the 25,331-image ISIC 2019 training release. The
aggregate collection includes images originating from HAM10000, BCN20000, and
MSK sources
\cite{tschandl2018ham10000,combalia2019bcn20000,codella2017isic}.

The original labels were mapped as follows:
NV, BKL, DF, and VASC formed the non-malignant class; MEL mapped to melanoma;
BCC mapped to BCC; and SCC mapped to SCC. Actinic keratosis was excluded
rather than forced into a malignant or non-malignant endpoint. Four images
belonging to label-conflicting duplicate components were also excluded,
leaving 24,460 eligible four-class images.

Images connected by a shared non-empty lesion identifier or identical
SHA-256 image hash were joined into split groups and assigned intact to the
training, validation, or internal-test partition using a diagnosis-stratified
70/15/15 procedure with seed 42. The audit found no cross-partition overlap in
split-group identifiers, available lesion identifiers, or exact hashes.
However, patient identifiers were unavailable and 2,084 images lacked lesion
identifiers. The procedure is therefore described as leakage-aware or
known-relation-disjoint, not fully patient-independent.

\begin{table}[t]
    \centering
    \caption{Primary four-class dataset distribution.}
    \label{tab:data_distribution}
    \footnotesize
    \setlength{\tabcolsep}{4.1pt}
    \begin{tabular}{lrrr}
        \toprule
        Class & Train & Validation & Test \\
        \midrule
        Non-malignant & 11,193 & 2,398 & 2,398 \\
        Melanoma      &  3,164 &   678 &   678 \\
        BCC           &  2,327 &   498 &   498 \\
        SCC           &    440 &    94 &    94 \\
        \midrule
        Total         & 17,124 & 3,668 & 3,668 \\
        \bottomrule
    \end{tabular}
\end{table}

All final four-class comparisons used the same locked 3,668-image internal-test
partition.

\subsection{Preprocessing}

All models used $224\times224$ RGB inputs and ImageNet normalisation.
Training transformations consisted of random resized cropping with scale
$0.85$--$1.00$ and aspect ratio $0.90$--$1.10$, horizontal and vertical
flips with probability 0.5, rotation within $\pm15^{\circ}$, and mild colour
jitter. Validation and test preprocessing deterministically resized the
shorter side to 256 pixels and applied a $224\times224$ centre crop. No
stochastic validation or test augmentation was used.

\subsection{Final Comparative Systems}

The final publication-facing four-class comparison uses two frozen systems.
The flat classifier maps every image directly to one of four endpoint labels:
non-malignant, melanoma, BCC, or SCC. The shared hierarchical model uses a
shared representation with task-specific prediction heads. Under deployed
hard routing, Task~1 first predicts non-malignant versus malignant. Samples
predicted malignant are passed to Task~2, which predicts melanoma, BCC, or
SCC; samples predicted non-malignant terminate at Task~1.

Task~3 predicts melanoma T categories and is analysed as a secondary task.
It is not part of the deployed four-class endpoint path. Shared Task~1--3
heads are also compared descriptively with their standalone task-model
counterparts to examine the trade-off between parameter sharing and
specialisation.

The final model identities, checkpoints, class order, preprocessing, routing
rules, and validation-based selection decisions were frozen before the
internal test was evaluated. Phase~06 manuscript integration introduces no
new training, inference rerun, checkpoint replacement, threshold search,
calibration, or model reselection.

\subsection{Oracle Routing Diagnostic}

To isolate the effect associated with upstream routing, an oracle-routing
condition retains the frozen downstream hierarchical predictors but replaces
the learned broad routing decision with the ground-truth broad category. For
non-malignant cases, the endpoint is therefore fixed as non-malignant; for
true malignant cases, the frozen Task~2 subtype prediction is used.

Oracle routing is diagnostic only. It provides an upper-bound-style estimate
of conditional hierarchical capability when Task~1 routing errors are removed,
but it is not an implementable or deployable inference policy and was not used
for any post-test tuning.

\subsection{Standalone T-Category Evidence}

The secondary melanoma T-category task uses the labels Tis, T1, T2, T3, and
T4. Its held-out test subset contains 127 cases. Rare-category support is
extremely limited, with T2, T3, and T4 supports of 7, 2, and 1, respectively.
Accordingly, Task~3 results are treated as support-limited descriptive evidence
rather than stable class-level performance estimates.

\subsection{Evaluation and Statistical Protocol}

Accuracy, balanced accuracy, macro-F1, weighted-F1, and per-class F1 were
reported for the final four-class comparison. Predictions from the flat model
and shared predicted-gate hierarchy were aligned by image identifier and
analysed using 10,000 ground-truth-class-stratified paired bootstrap
replicates with seed 42 and percentile 95\% confidence intervals for metric
differences.

Exact two-sided McNemar testing was applied to paired correctness outcomes.
McNemar's test is therefore interpreted as evidence about correctness
asymmetry and the accuracy comparison; it is not used as a significance test
for macro-F1.

The internal test was consumed once under the frozen protocol. Its results
were not used for threshold tuning, checkpoint selection, architecture
selection, preprocessing changes, calibration, post-hoc optimisation,
retraining, or model reselection. Validation-to-test comparisons are reported
only as directional within-study consistency and were not used to modify the
frozen systems.
